\documentclass[11pt]{article}
\usepackage[utf8]{inputenc}
\usepackage[T1]{fontenc}
\usepackage{lmodern}
\usepackage[spanish]{babel}
\usepackage[margin=2.5cm]{geometry}
\usepackage{amsmath,amssymb}
\usepackage{parskip}
\usepackage{xcolor}
\usepackage[colorlinks=true,linkcolor=black,urlcolor=blue,citecolor=black]{hyperref}

\title{\textbf{El arte de las matemáticas:\\de la proporción áurea al arte generativo}}
\author{Antonio Soria\\ \small{Divulgación Científica — purpplealien.github.io}}
\date{Julio 2026}

\begin{document}
\maketitle

\begin{abstract}
Arte y matemáticas suelen presentarse como opuestos: emoción contra rigor, intuición contra
demostración. La historia cuenta otra cosa: la perspectiva nació de la geometría, la armonía
musical de las proporciones, y hoy los algoritmos pintan. Este artículo recorre esa relación
en cuatro estaciones --- la proporción áurea, los fractales, las teselaciones de Escher y el
arte generativo --- y termina donde trabajo yo: escribiendo código que dibuja.
\end{abstract}

\section{La proporción áurea: mito y matemática}

Pocas constantes cargan tanta leyenda como
\[
\varphi = \frac{1 + \sqrt{5}}{2} \approx 1.6180339887\ldots
\]
Su definición es pura elegancia: es la única manera de dividir un segmento tal que el todo es
a la parte mayor como la parte mayor es a la menor, lo que conduce a la ecuación
$\varphi^2 = \varphi + 1$. De ahí brotan propiedades deliciosas: $\varphi = 1 + 1/\varphi$,
su expansión en fracción continua es la más simple posible, $[1;1,1,1,\ldots]$, y eso la
convierte --- en un sentido técnico preciso --- en \emph{el número más difícil de aproximar
con fracciones}: el más irracional de los irracionales.

Está íntimamente ligada a la sucesión de Fibonacci $1, 1, 2, 3, 5, 8, 13, 21, \ldots$, pues
el cociente de términos consecutivos converge a $\varphi$:
\[
\lim_{n \to \infty} \frac{F_{n+1}}{F_n} = \varphi.
\]

Conviene la honestidad divulgativa: buena parte de las apariciones de $\varphi$ en el arte
clásico (el Partenón, La Gioconda) son dudosas o directamente mitología moderna --- medir un
rectángulo aproximado en una fachada no demuestra intención. Donde $\varphi$ sí aparece con
rigor demoledor es en la naturaleza, y por una razón funcional que veremos al final: la
disposición de las semillas del girasol, las escamas de la piña, las hojas alrededor del
tallo.

\section{Fractales: la geometría de lo rugoso}

La geometría clásica describe círculos y rectas, pero --- como escribió Benoît Mandelbrot ---
``las nubes no son esferas, las montañas no son conos''. Su respuesta (1975) fueron los
\emph{fractales}: objetos autosemejantes cuyo detalle no se agota al hacer zoom, con
dimensiones que desafían la intuición. El copo de nieve de Koch tiene dimensión
\[
D = \frac{\log 4}{\log 3} \approx 1.2619:
\]
más que una línea, menos que una superficie; un perímetro infinito encerrando un área finita.

Los fractales resultaron ser, además, profundamente \emph{bellos}. El conjunto de Mandelbrot
--- definido por la iteración trivial $z_{n+1} = z_n^2 + c$ sobre números complejos --- generó
en los años ochenta una estética completa de pósters, portadas de discos y salvapantallas.
Y hay evidencia empírica de que la preferencia no es casual: los estudios sobre las pinturas
de goteo de Jackson Pollock encontraron en ellas dimensiones fractales consistentes
($D \approx 1.3$--$1.7$ según la época), y experimentos de percepción sugieren que los
patrones con $D \approx 1.3$--$1.5$ --- el rango de nubes, costas y copas de árboles ---
resultan especialmente placenteros y hasta reductores de estrés. Pollock, sin saberlo,
pintaba la estadística del paisaje.

\section{Escher: el arte de agotar el plano}

Maurits Cornelis Escher no sabía matemáticas formales, y sin embargo su obra dialoga con dos
teorías profundas. Sus teselaciones --- lagartos, aves y peces que llenan el plano sin huecos
--- exploran sistemáticamente los \emph{grupos cristalográficos}: las 17 formas esencialmente
distintas de repetir un motivo en el plano, clasificadas por la teoría de grupos en el siglo
XIX. Escher las redescubrió a mano, copiando y transformando los mosaicos de la Alhambra.

Su serie \emph{Límite circular} fue más lejos: tras conocer al geómetra H.~S.~M.~Coxeter,
Escher llenó discos con figuras que se encogen hacia el borde siguiendo con precisión el
modelo de Poincaré de la \emph{geometría hiperbólica} --- un universo donde por un punto pasan
infinitas paralelas y que cabe, conformemente, dentro de un círculo. Coxeter analizó después
los grabados y confirmó que la construcción era matemáticamente exacta. Pocas colaboraciones
ilustran mejor la tesis de este artículo.

\section{Arte generativo: cuando el pincel es un algoritmo}

Desde los años sesenta --- con pioneros como Vera Molnár, Frieder Nake y Georg Nees --- existe
arte cuyo material no es el óleo sino el \emph{proceso}: reglas, aleatoriedad controlada y una
máquina que ejecuta. Hoy ese campo tiene revistas académicas (\emph{Leonardo}, del MIT),
congresos dedicados (\emph{Bridges}, sobre matemáticas y arte) y herramientas libres
(Processing, p5.js) usadas por miles de artistas.

Lo fascinante es cuánta belleza brota de fórmulas mínimas. Una \emph{rosa polar} es solo
$r = \cos(k\theta)$; las curvas de Lissajous, dos osciladores desfasados; un atractor extraño,
tres ecuaciones diferenciales iteradas millones de veces. En mi sitio web mantengo un motor de
arte generativo con ocho piezas de esta familia --- rosas de Maurer, espirógrafos, la mariposa
de Fay, los atractores de Lorenz y De Jong, y una filotaxis --- dibujadas en vivo por el
navegador, con los parámetros a la vista para que cualquiera experimente.

\subsection{La filotaxis: el ejemplo perfecto}

La imagen que acompaña este artículo en mi página es una \emph{filotaxis}: el patrón de las
semillas del girasol. Su receta completa cabe en dos líneas. La semilla $n$-ésima se coloca en
\[
r = c\sqrt{n}, \qquad \theta = n \cdot 137.507\ldots^{\circ},
\]
donde el ángulo mágico es el \emph{ángulo áureo}, $360^{\circ}/\varphi^2$. Aquí cierra el
círculo del artículo: como $\varphi$ es ``el más irracional'' de los números, girar por el
ángulo áureo es la manera óptima de \emph{nunca} alinear dos semillas en el mismo rayo ---
cada semilla nueva cae siempre en el hueco más grande disponible. Las plantas no conocen
fracciones continuas; la selección natural optimizó el empaquetamiento y la matemática estaba
ahí, esperando ser leída. El resultado, además de óptimo, es hermoso: espirales entrelazadas
cuyos conteos son números de Fibonacci consecutivos.

\section{Conclusión}

La proporción áurea ordena girasoles, los fractales explican por qué Pollock relaja, la teoría
de grupos cataloga los mosaicos de la Alhambra y dos líneas de trigonometría dibujan un jardín
digital. La frontera entre arte y matemáticas nunca existió: existen, apenas, dos lenguajes
distintos para la misma obsesión humana --- encontrar y crear patrones. Y hoy, con un
navegador y cincuenta líneas de código, esa frontera inexistente es también un territorio
que cualquiera puede habitar.

\section*{Para profundizar}

\begin{itemize}
  \item Bridges Organization (congreso anual de matemáticas y arte): \url{https://www.bridgesmathart.org/}
  \item Revista \emph{Leonardo}, MIT Press: \url{https://leonardo.info/}
  \item R.~P.~Taylor et al., \emph{Fractal analysis of Pollock's drip paintings}, Nature 399 (1999).
  \item Mi motor de arte generativo en vivo: \url{https://purpplealien.github.io/algoritmos.html}
\end{itemize}

\end{document}
